% Owner: member A.
\subsection{(a)(1) Problem definition, challenges and solution types}

\paragraph{Problem definition.}
Sentiment analysis, also called opinion mining, is the computational study of the opinions, sentiments and emotions that people express in text~\cite{pang2008opinion,liu2012sentiment}.
Liu~\cite{liu2012sentiment} formalises an opinion as a quintuple $(e, a, s, h, t)$: an entity $e$, one of its aspects $a$, the sentiment $s$ towards that aspect, the opinion holder $h$ and the time $t$.
This project is the simplest instance of that definition: \emph{document-level binary polarity classification}.
The entity is the purchased product and the holder is the reviewer, both implicit, and each review receives one label for the whole document.
Given a review $x$, we learn a classifier $f(x) \approx P(y{=}1 \mid x)$, where $y=1$ is positive and $y=0$ is negative, and predict $\hat y = \mathbb{1}[f(x) \ge \tau]$ for a decision threshold $\tau$.
Neutral or irrelevant reviews are not modelled.

\paragraph{Granularity.}
Sentiment can be analysed at three levels of granularity:
\begin{itemize}
  \item \emph{Document level}: one label per review, as in this project~\cite{pang2002thumbs}.
  \item \emph{Sentence level}: usually a subjectivity filter followed by polarity classification.
  \item \emph{Aspect level}: aspect-based sentiment analysis (ABSA) extracts aspects such as ``fit'' or ``fabric'' and assigns a polarity to each~\cite{hu2004mining,pontiki2014semeval}.
\end{itemize}
Related tasks use finer label sets, such as the five-point scale of the Stanford Sentiment Treebank~\cite{socher2013recursive}, or emotion categories.

\paragraph{Challenges.}
The main difficulties, each illustrated with what we observe in our data:
\begin{enumerate}
  \item \textbf{Negation and contrast.} Words such as ``not'', ``never'' and ``but'' change the polarity of the words around them. In ``Not bad for just a shirt'' (from our test set), ``bad'' signals a negative review, yet the review is mildly positive. A bag-of-words model loses the word order needed to resolve this.
  \item \textbf{Implicit sentiment, sarcasm and irony.} Many reviews contain no explicit opinion word (``the seam came apart after one wash''), or say the opposite of what they mean.
  \item \textbf{Mixed polarity.} A review may praise one aspect and criticise another. The document label then depends on which aspect the reviewer weighs more.
  \item \textbf{Domain dependence.} The polarity of a word depends on the domain. ``Small'' is negative for a hotel room and often positive for a phone, and ``unpredictable'' is positive for a film plot but negative for a car's steering. Classifiers trained on one product category therefore degrade on another~\cite{turney2002thumbs,blitzer2007biographies}.
  \item \textbf{Class imbalance.} Product reviews are mostly positive: 85.4\% of our training reviews are positive and 14.6\% negative. Predicting ``positive'' for every review already gives 85.4\% accuracy, so we report Macro-F1 and the F1 of the negative class alongside accuracy.
  \item \textbf{Noisy and informal text.} Reviews contain typos, slang, emoticons and very different lengths. Our training reviews have a median of 42 words, but the longest has 1{,}066. When star ratings are used as labels, some labels disagree with the text, which is the problem studied in (i).
\end{enumerate}

\paragraph{Solution types under different settings.}
Table~\ref{tab:a1-settings} organises common approaches along the three axes named in the assignment.

\emph{Supervised versus unsupervised.}
Supervised methods learn from labelled reviews. They range from Naive Bayes, maximum-entropy and SVM classifiers on bag-of-words features~\cite{pang2002thumbs}, and their strong NB-SVM variant~\cite{wang2012baselines}, through convolutional and recurrent networks over word embeddings~\cite{kim2014convolutional,socher2013recursive,zhang2018deep}, to fine-tuning pre-trained Transformers such as BERT and RoBERTa~\cite{devlin2019bert,liu2019roberta}.
Unsupervised, lexicon-based methods need no labels:
\begin{itemize}
  \item scoring phrases by their association with seed words such as ``excellent'' and ``poor''~\cite{turney2002thumbs};
  \item looking words up in a sentiment lexicon such as SentiWordNet~\cite{baccianella2010sentiwordnet};
  \item rule-based scoring with VADER~\cite{hutto2014vader}.
\end{itemize}
They transfer easily across domains but are usually less accurate than supervised models.
Between the two extremes lie weakly supervised methods. Distant supervision uses emoticons or star ratings as noisy labels~\cite{go2009twitter}, and semi-supervised methods use self-training on unlabelled reviews. Zero-shot and few-shot prompting of large language models is competitive on simple binary polarity, but still falls behind fine-tuned models on fine-grained and structured sentiment tasks~\cite{zhang2023sentiment}.

\emph{Closed-set versus open-set.}
In the closed-set setting every test review belongs to one of the training classes, here positive or negative.
In the open-set setting, test inputs may come from classes never seen in training: neutral reviews, spam, text that is not a review, or another language.
A closed-set classifier still forces each of these into ``positive'' or ``negative''.
Open-set methods add a rejection option. DOC replaces the softmax with one-vs-rest sigmoids and rejects an input when no class is confident enough~\cite{shu2017doc}. Confidence-based detection flags inputs whose maximum softmax probability is low~\cite{hendrycks2017baseline}.

\emph{Without versus with domain shift.}
When training and test data come from the same distribution, standard supervised learning applies.
Under domain shift, labelled data from a source domain (e.g.\ books) must serve an unlabelled target domain (e.g.\ kitchen appliances). Approaches include:
\begin{itemize}
  \item structural correspondence learning, which aligns domain-specific words through shared ``pivot'' features~\cite{blitzer2007biographies};
  \item stacked denoising autoencoders that learn domain-invariant representations~\cite{glorot2011domain};
  \item domain-adversarial training (DANN)~\cite{ganin2016domain};
  \item continued pre-training of the language model on target-domain text~\cite{gururangan2020dont}.
\end{itemize}
These methods are the basis of our answer to (h).

\begin{table}[t]
\centering
\caption{Solution types for sentiment classification under different settings. Our project is marked with $\star$.}
\label{tab:a1-settings}
\small
\begin{tabularx}{\textwidth}{l Y Y}
\toprule
Setting & Typical methods & Strengths / limitations \\
\midrule
Supervised $\star$ & BoW + NB/LR/SVM~\cite{pang2002thumbs,wang2012baselines}; CNN/RNN over embeddings~\cite{kim2014convolutional,socher2013recursive}; fine-tuned Transformers~\cite{devlin2019bert,liu2019roberta} & Most accurate when labelled in-domain data exist; needs labels and is domain-specific. \\
Unsupervised & Seed-word association~\cite{turney2002thumbs}; lexicons~\cite{baccianella2010sentiwordnet}; rules (VADER)~\cite{hutto2014vader} & No labels needed and interpretable; misses context, domain-specific terms and sarcasm. \\
Weak / semi-supervised & Distant labels from emoticons or ratings~\cite{go2009twitter}; self-training; LLM prompting~\cite{zhang2023sentiment} & Cheap labels at scale; labels are noisy (see (i)). \\
\midrule
Closed-set $\star$ & Softmax over a fixed label set & Simple; inputs from unseen classes are forced into a known class. \\
Open-set & One-vs-rest rejection (DOC)~\cite{shu2017doc}; confidence thresholds~\cite{hendrycks2017baseline} & Can reject neutral or irrelevant inputs; needs extra thresholds. \\
\midrule
No domain shift $\star$ & Standard train/validation/test protocol & Validation accuracy predicts test accuracy. \\
Domain shift & SCL pivots~\cite{blitzer2007biographies}; denoising autoencoders~\cite{glorot2011domain}; DANN~\cite{ganin2016domain}; domain-adaptive pre-training~\cite{gururangan2020dont} & Uses unlabelled target text; gains depend on how far apart the domains are. \\
\bottomrule
\end{tabularx}
\end{table}

\paragraph{Our setting.}
The project is \emph{supervised, closed-set and in-domain}.
Training and test reviews come from the same product-review distribution; 101 of the 1{,}851 test reviews even appear verbatim in the training set.
The data are imbalanced, and the labels may contain some noise.
We therefore build supervised classifiers. A TF-IDF linear model serves as the baseline, and sequence models and pre-trained language models serve as improvements; (a)(4) motivates this choice.
We use threshold tuning and Macro-F1 to handle the imbalance.
